\documentclass[12pt,a4paper]{article}
\usepackage[utf8]{inputenc}
\usepackage{amsmath,amssymb,amsthm}
\usepackage{hyperref}
\usepackage{booktabs}
\usepackage{graphicx}
\usepackage{natbib}

\newtheorem{theorem}{Theorem}
\newtheorem{definition}{Definition}
\newtheorem{corollary}{Corollary}

\newcommand{\emlstar}{\mathrm{eml}^{\star}}
\newcommand{\eml}{\mathrm{eml}}
\newcommand{\conj}[1]{\overline{#1}}

\title{Detecting anti-holomorphic dynamics via symbolic regression\\
with the $\emlstar$ operator}

\author{Anthony Monnerot\\
Independent researcher, France\\
\texttt{github.com/antparis/eml\_star}}

\date{May 2026}

\begin{document}
\maketitle

\begin{abstract}
We introduce $\emlstar(x,y) = \exp(\conj{x}) - \ln(\conj{y})$,
the anti-holomorphic extension of the EML operator
$\eml(x,y) = \exp(x) - \ln(y)$ of Odrzywołek (2026).
We prove that $\conj{z} = 1 - \emlstar(0, \eml(z,1))$
(Theorem~\ref{thm:conj}), providing a purely algebraic
representation of complex conjugation within the EML framework.
Combined with symbolic regression (PySR), this yields an automatic
criterion for classifying complex dynamical systems:
$\emlstar$ appears in the discovered equation when the underlying
dynamics is anti-holomorphic and does not appear when it is
holomorphic, across all tested systems.
We validate this on three domains: (i) fractal iteration maps,
where PySR distinguishes the Mandelbrot set ($z^2+c$, holomorphic,
$\emlstar$-free) from the Tricorn ($\conj{z}^2+c$, anti-holomorphic,
$\emlstar$-present) at machine precision ($\mathrm{MSE} \sim 10^{-32}$);
(ii) quantum state evolution, where $\emlstar$ correctly identifies
dissipative systems (7/7 classification);
(iii) gravitational lensing ellipticities from KiDS-1000,
confirming $\emlstar$ detection on 5\,000 real astrophysical
measurements.
A negative control on the weak gravitational shear profile around
KiDS-1000 lenses (268\,722 source-lens pairs) shows that $\emlstar$
does not appear at low complexity when the underlying physics is
holomorphic (standard general relativity), confirming specificity.
\end{abstract}


% =====================================================================
\section{Introduction}
\label{sec:intro}

A fundamental distinction in complex dynamics is between holomorphic
maps, which preserve the complex structure, and anti-holomorphic maps,
which reverse it through complex conjugation.
This distinction governs whether a system conserves information
(unitary evolution in quantum mechanics, conformal maps in fluid
dynamics) or dissipates it (damped oscillations, measurement collapse).

Despite its importance, no automatic method existed to determine
from data alone whether a governing equation is holomorphic or
anti-holomorphic.
Standard symbolic regression tools~\citep{cranmer2023pysr} search over
real-valued operators and cannot distinguish $z^2$ from $\conj{z}^2$
since these coincide on the real line.

Odrzywołek~\citep{odrzywołek2026eml} showed that a single binary
operator $\eml(x,y) = \exp(x) - \ln(y)$ generates all elementary
functions through uniform binary trees.
We extend this framework with $\emlstar$, the anti-holomorphic
companion of $\eml$, and show that it provides the missing piece:
an intrinsic marker of anti-holomorphic structure that appears
automatically in symbolic regression when the underlying dynamics
requires complex conjugation.


% =====================================================================
\section{The $\emlstar$ operator}
\label{sec:emlstar}

\begin{definition}[EML operator]
Following~\citet{odrzywołek2026eml}, the EML operator is:
\begin{equation}
\eml(x,y) = \exp(x) - \ln(y), \quad y \neq 0.
\end{equation}
\end{definition}

\begin{definition}[$\emlstar$ operator]
The anti-holomorphic extension of EML is:
\begin{equation}
\emlstar(x,y) = \exp(\conj{x}) - \ln(\conj{y}), \quad y \neq 0.
\label{eq:emlstar}
\end{equation}
\end{definition}

\noindent
Note that $\emlstar$ coincides with $\eml$ on $\mathbb{R}$,
since $\conj{x} = x$ for $x \in \mathbb{R}$.
This is not a defect but a feature:
$\emlstar$ can only be distinguished from $\eml$ on data with
genuine complex structure.

\begin{theorem}[Conjugation via $\emlstar$]
\label{thm:conj}
For all $z \in \mathbb{C}$ with $|\mathrm{Im}(z)| < \pi$:
\begin{equation}
\conj{z} = 1 - \emlstar\!\left(0,\; \eml(z, 1)\right).
\label{eq:conj}
\end{equation}
\end{theorem}

\begin{proof}
We compute stepwise:
\begin{align}
\eml(z, 1) &= \exp(z) - \ln(1) = \exp(z), \\
\emlstar(0, \exp(z)) &= \exp(\conj{0}) - \ln(\conj{\exp(z)}) \\
  &= 1 - \ln(\exp(\conj{z})) \\
  &= 1 - \conj{z},
\end{align}
where the last step uses $\ln(\exp(w)) = w$ on the principal branch
($|\mathrm{Im}(w)| < \pi$, satisfied since $\mathrm{Im}(\conj{z}) = -\mathrm{Im}(z)$).
Therefore $1 - \emlstar(0, \eml(z,1)) = 1 - (1 - \conj{z}) = \conj{z}$.
\end{proof}

\begin{corollary}[Modulus squared]
$|z|^2 = z \cdot \conj{z} = z \cdot [1 - \emlstar(0, \eml(z,1))]$.
\end{corollary}

\begin{corollary}[Density]
\label{cor:density}
The algebra $\mathcal{A}$ generated by $\{\eml, \emlstar, 1\}$ is dense in
$C(K, \mathbb{C})$ for any compact $K \subset \mathbb{C}$ contained
in a horizontal strip of height less than $2\pi$,
by the Stone--Weierstrass theorem for complex algebras.
The algebra separates points (since $\eml(z,1) = \exp(z)$ is injective
on such strips), contains the constants, and is self-adjoint
(closed under conjugation, since $\emlstar$ provides $\conj{z}$
via Theorem~\ref{thm:conj}).
\end{corollary}

% =====================================================================
\section{A Re-based decomposition of $\emlstar$}
\label{sec:re-decomposition}

The $\emlstar$ operator was introduced in Section~\ref{sec:emlstar}
as the natural anti-holomorphic companion of $\eml$.
A symbolic regression experiment using PySR was run with $z$ and $w$
as independent complex inputs sampled from the square
$\{|\mathrm{Re}|, |\mathrm{Im}| \leq \pi - 0.05\}$ with $|w| \geq 0.1$,
and the toolbox restricted to
$\{+, -, \times, /, \sin, \cos, \exp, \log,
  \mathrm{my\_real}, \mathrm{my\_imag}\}$
(no \texttt{my\_conj}, no $\eml$, no $\emlstar$).
On 2\,000 training points, PySR rediscovered the following
identity at $\mathrm{MSE} = 4.22 \times 10^{-31}$
(machine precision for complex128):

\begin{equation}
\boxed{\,\emlstar(z, w) = \mathrm{Re}\!\left(2\exp(z) - \ln(w^2)\right)
+ \ln(w) - \exp(z)\,}
\label{eq:re-decomp}
\end{equation}

An independent verification on $N = 10\,000$ points
(different seed) confirmed $\mathrm{MSE} = 3.17 \times 10^{-31}$
and maximum absolute error $7.11 \times 10^{-15}$.

\paragraph{Algebraic proof.}
For any $u \in \mathbb{C}$, the elementary identity
$\conj{u} = 2\,\mathrm{Re}(u) - u$ holds.
Applying it to $u = \exp(z)$ and using
$\conj{\exp(z)} = \exp(\conj{z})$ yields

\begin{equation}
\exp(\conj{z}) = 2\,\mathrm{Re}(\exp(z)) - \exp(z).
\tag{Lemma A}
\label{eq:lemma-A}
\end{equation}

Applying it to $u = \ln(w)$, and using
$\conj{\ln(w)} = \ln(\conj{w})$ on the principal branch
($\arg(w) \in (-\pi, \pi)$), yields

\begin{equation}
\ln(\conj{w}) = 2\,\mathrm{Re}(\ln(w)) - \ln(w).
\tag{Lemma B}
\label{eq:lemma-B}
\end{equation}

Substituting Lemma A and Lemma B into the definition
$\emlstar(z, w) = \exp(\conj{z}) - \ln(\conj{w})$
gives equation~\eqref{eq:re-decomp}. $\square$

\paragraph{Structural consequence.}
Equation~\eqref{eq:re-decomp} expresses $\emlstar$ using only
$\mathrm{Re}$, $\exp$, $\ln$, and arithmetic operations.
Since $\{\eml, 1\}$ generates $\exp$ and $\ln$
(Odrzywołek, 2026), it follows that
$\{\eml, \mathrm{Re}, 1\}$ generates $\emlstar$.
Combined with the reverse direction
($\mathrm{Re}$ is recoverable from $\{\eml, \emlstar, 1\}$
via $\mathrm{Re}(z) = (z + \conj{z})/2$ and Theorem~\ref{thm:conj}),
we obtain:

\begin{theorem}[Equivalence of minimal non-holomorphic extensions]
\label{thm:equivalence}
The operator sets $\{\eml, \emlstar, 1\}$ and
$\{\eml, \mathrm{Re}, 1\}$ have identical expressive power
over $C(K, \mathbb{C})$ for every compact
$K \subset \mathbb{C}$ contained in a horizontal strip
of height less than $2\pi$ and disjoint from $\{w = 0\}$.
\end{theorem}

\begin{corollary}[Generalized completeness]
\label{cor:generalized}
Let $\phi$ be any non-holomorphic primitive in
$\{\mathrm{Re}, \mathrm{Im}, \conj{\cdot}, \emlstar\}$.
The set $\{\eml, \phi, 1\}$ is dense in $C(K, \mathbb{C})$
for every compact $K$ satisfying the branch condition above.
\end{corollary}

\paragraph{Discussion.}
Theorem~\ref{thm:equivalence} relativizes the apparent uniqueness
of $\emlstar$ as the natural anti-holomorphic extension of $\eml$.
Any single non-holomorphic primitive added to $\{\eml, 1\}$
achieves $C(K, \mathbb{C})$-density.
The $\emlstar$ formulation retains structural appeal
(it preserves the binary, exp-minus-log signature of $\eml$),
while the $\mathrm{Re}$ formulation is more minimal in the
standard analytic toolkit.

The empirical rediscovery of equation~\eqref{eq:re-decomp}
by an unconstrained symbolic regression search demonstrates
that the $\mathrm{Re}$-based decomposition is not merely
formally valid but \emph{algorithmically accessible}:
a generic evolutionary search over elementary operators
converges to it within minutes when given enough expressive flexibility.
Implementation and reproducibility scripts are available at
\url{https://github.com/antparis/oxieml-star}
(commit \texttt{f03a46e}).



% =====================================================================
\section{Method}
\label{sec:method}

\subsection{Symbolic regression with complex operators}

We use PySR~\citep{cranmer2023pysr} (v1.5.10) with Julia
backend, configured for complex arithmetic:

\begin{itemize}
\item \textbf{Binary operators:} $+, -, \times, \div, \eml, \emlstar$
\item \textbf{Unary operators:} $\cos, \sin, \exp, \log,
  \mathrm{my\_conj}(z) = \conj{z},\;
  \mathrm{my\_real}(z) = \mathrm{Re}(z),\;
  \mathrm{my\_imag}(z) = \mathrm{Im}(z),\;
  \mathrm{my\_abs2}(z) = z\conj{z}$
\item \textbf{Precision:} \texttt{Complex\{Float64\}}
\item \textbf{Reproducibility:} \texttt{deterministic=True,
  parallelism="serial", random\_state=42}
\end{itemize}

\subsection{Detection criterion}

A system is classified as \emph{anti-holomorphic} if the best
equation found by PySR contains any of: \texttt{emlstar},
\texttt{my\_conj}, \texttt{my\_abs2}, or \texttt{my\_imag}.
These all involve complex conjugation.

\subsection{Why $\emlstar$ is necessary}

Standard symbolic regression over $\{+,-,\times,\div,\exp,\sin,\cos,\log\}$
cannot distinguish $f(z) = z^2$ from $g(z) = \conj{z}^2$, because
all standard operators are holomorphic.
Adding $\emlstar$ to the operator set is the minimal extension that
makes anti-holomorphic functions expressible.
By Corollary~\ref{cor:density}, $\{\eml, \emlstar\}$ suffices
to approximate any continuous complex function.


% =====================================================================
\section{Results}
\label{sec:results}

\subsection{Fractal iteration maps}
\label{sec:fractals}

We generate 200 random points $z \in [-1,1]^2 \subset \mathbb{C}$
(seed 42) and compute one iteration of three fractal maps with
$c = -0.7 + 0.27015i$:

\begin{table}[h]
\centering
\caption{Fractal maps: PySR discovers the exact iteration rule
and correctly identifies the holomorphic/anti-holomorphic nature.}
\label{tab:fractals}
\begin{tabular}{lllcr}
\toprule
Fractal & True rule & PySR equation & $\emlstar$? & MSE \\
\midrule
Mandelbrot & $z^2 + c$ & $(x_0 \cdot x_0) - (0.7 - 0.27i)$ & No & $1.86 \times 10^{-32}$ \\
Tricorn & $\conj{z}^2 + c$ & $\mathrm{my\_conj}((x_0 \cdot x_0) - (0.7 + 0.27i))$ & \textbf{Yes} & $2.02 \times 10^{-32}$ \\
Burning Ship & $(|\mathrm{Re}|+i|\mathrm{Im}|)^2+c$ & approximation & \textbf{Yes} & $1.39 \times 10^{-2}$ \\
\bottomrule
\end{tabular}
\end{table}

The Mandelbrot and Tricorn results are at machine precision.
PySR recovers the exact algebraic form and uses $\mathrm{my\_conj}$
(equivalent to $\emlstar$ via Theorem~\ref{thm:conj}) for the Tricorn
but not for the Mandelbrot.
This is the cleanest possible test: data is natively complex,
no encoding artefact is possible, and the only difference between
the two maps is one conjugation.


\subsection{Quantum state evolution}
\label{sec:evolution}

We generate time series $\psi(t) \to \psi(t+dt)$ for three
quantum systems ($\omega = 1.0$, $dt = 0.1$, 300 time steps):

\begin{table}[h]
\centering
\caption{Quantum evolution: PySR finds the exact evolution operator
and identifies dissipative systems via $\emlstar$.}
\label{tab:evolution}
\begin{tabular}{llccl}
\toprule
System & $|U|$ & $\emlstar$? & MSE & Physics \\
\midrule
Larmor precession & 1.0000 & No & $5.24 \times 10^{-32}$ & Unitary \\
Damped oscillation & 0.9900 & \textbf{Yes} & $1.69 \times 10^{-32}$ & Dissipative \\
Rabi oscillation & 1.0000 & No & $2.79 \times 10^{-31}$ & Unitary \\
\bottomrule
\end{tabular}
\end{table}

\begin{table}[h]
\centering
\caption{Multi-domain physics: additional natively complex systems.}
\label{tab:multidomain}
\begin{tabular}{llcr}
\toprule
System & Observable & $\emlstar$? & MSE \\
\midrule
Impedance (RLC circuit) & $|Z(\omega)|^2$ & \textbf{Yes} & $2.48 \times 10^{-31}$ \\
Bell state (Born rule) & $|\psi|^2$ & \textbf{Yes} & $1.74 \times 10^{-33}$ \\
ECG analytic signal & $|z(t)|^2$ & \textbf{Yes} & $5.58 \times 10^{-34}$ \\
EM plane wave & $|E+iB|^2 = 1$ (const.) & No & $1.97 \times 10^{-33}$ \\
\bottomrule
\end{tabular}
\end{table}

All seven systems are classified correctly.
$\emlstar$ appears when the physics requires non-trivial
conjugation: dissipation ($|U| < 1$) or measurement ($|z|^2$
with varying $z$). For the EM plane wave, $|E+iB|^2 = 1$
is constant (the wave is on the unit circle), so PySR
discovers the constant directly without conjugation.
It does not appear for unitary or conservative systems.


\subsection{Astrophysical validation: KiDS-1000}
\label{sec:kids}

As a first application to real astrophysical data, we use 5\,000
galaxy ellipticities from the KiDS-DR4 DEIMOS catalog
\citep{giblin2021kids1000}: each galaxy has a measured shape
$\gamma = e_1 + i\,e_2 \in \mathbb{C}$.
We test whether PySR recovers $|\gamma|^2 = \gamma \cdot \conj{\gamma}$:

\begin{equation}
\text{PySR finds: } \mathrm{my\_abs2}(x_0), \quad
\mathrm{MSE} = 2.54 \times 10^{-34}, \quad \emlstar = \textbf{Yes}.
\end{equation}

This confirms that $\emlstar$ operates correctly on real
astrophysical measurements with native complex structure.
While the test is a sanity check ($|\gamma|^2$ is known to
require conjugation), it is the first application of $\emlstar$
to observational data.


% =====================================================================
\section{Negative control: weak gravitational shear profile}
\label{sec:lensing-control}

The previous tests show $\emlstar$ appearing on systems where
anti-holomorphic structure is expected from theory.
A stronger test of specificity is to apply the same machinery to
data where anti-holomorphic structure is \emph{not} expected and
verify that $\emlstar$ does not appear.

We use the KiDS-1000 SOM-gold catalog~\citep{giblin2021kids1000}, comprising
21\,262\,011 source galaxies with reliable shape and photometric
redshift measurements over 1006 deg$^2$, and the 268 high-quality
strong-lens candidates of~\citet{petrillo2019linKS, li2020kids, li2021kids}.
After standard selection cuts
($\mathtt{fitclass}=0$, $\mathtt{MASK}=0$, $\mathtt{weight}>0$,
$\mathtt{Z\_B}>0.3$) we retain 17\,450\,042 sources.
Pairing every source within $10$ arcmin of each lens yields
293\,880 (lens, source) pairs spread across 181/268 lenses
with at least one matched source. Of these, 268\,722 pairs fall within the analysis range $r \in [0.5, 10]$ arcmin used for stacking.

For each pair, we compute the lens--source angular separation
$\Delta z = (\mathrm{RA}_s - \mathrm{RA}_\ell) + i(\mathrm{Dec}_s - \mathrm{Dec}_\ell)$
and the rotated source ellipticity
$\gamma_\mathrm{rot} = \gamma \exp(-2i\varphi)$,
where $\gamma = e_1 + ie_2$ and $\varphi = \arg(\Delta z)$.
We stack pairs in 20 logarithmic bins of $r = |\Delta z|$
between 0.5 and 10 arcmin, using \texttt{lensfit} inverse-variance
weights:
\begin{equation}
\langle\gamma_\mathrm{rot}\rangle(r)
\;=\;
\frac{\sum_k w_k\,\gamma_{\mathrm{rot},k}}{\sum_k w_k}.
\label{eq:stacking}
\end{equation}
All 20 bins have effective stacking weight $N_\mathrm{eff} > 100$
and are retained.

\paragraph{Anti-circularity.}
The rotation factor $\exp(-2i\varphi) = (\overline{\Delta z}/|\Delta z|)^2$
contains an implicit conjugation.
If $\Delta z$ or $\varphi$ were passed to PySR, the system could
trivially reconstruct $\gamma_\mathrm{rot}$ via complex conjugation,
giving a circular result.
To prevent this, PySR receives only $r = |\Delta z|$ as a single
real-valued input (cast to \texttt{complex128} with zero imaginary
part), with $\langle\gamma_\mathrm{rot}\rangle(r)$ as the complex
target.
This is verified in code by asserting $\mathrm{Im}(X) = 0$ before
fitting.

\paragraph{Result.}
PySR with the full $\emlstar$-enabled toolbox is run on the 20-bin
stacked profile.
The Pareto front shows a clear holomorphic plateau at low complexity:
the formula $(c_1 + ic_2) / (x_0 + c_3)$ at complexity 5 reaches
$\mathrm{MSE} = 5.7 \times 10^{-5}$, a factor 2.4 improvement over
the constant baseline.
This is the expected $1/r$ shear scaling from standard general
relativity.
Complexity-11 reaches $\mathrm{MSE} = 3.9 \times 10^{-5}$ with a
purely holomorphic formula and \emph{beats} the only
anti-holomorphic candidate at complexity 10
(MSE $4.5 \times 10^{-5}$, score 0.16).
The $\emlstar$ operator does appear at complexity $\geq 14$, but
with a marginal score per added node ($\lesssim 0.03$) and within
the regime where PySR can interpolate noise on 20 binned points.

\paragraph{Interpretation.}
This is the contrapositive of the positive tests of
Section~\ref{sec:results}.
On a real astrophysical signal that is well described by
holomorphic general relativity, $\emlstar$ is \emph{not} structurally
required: PySR prefers simpler holomorphic formulas at low complexity
and only resorts to $\emlstar$ for cosmetic high-complexity fits.
This is the behavior expected of a specific anti-holomorphic
detector, not a free-floating fit parameter.



% =====================================================================
\section{Discussion}
\label{sec:discussion}

\subsection{What $\emlstar$ detects}

Across all validated tests, $\emlstar$ appears when
the dynamics involves complex conjugation and does not appear
when it does not:
\begin{itemize}
\item Anti-holomorphic maps ($\conj{z}^2 + c$)
\item Dissipative evolution ($|U| < 1$, requiring $|z|^2 = z\conj{z}$)
\item Born rule measurements ($P = |\psi|^2$)
\end{itemize}
It does not appear for holomorphic, unitary, or energy-conserving systems.

\subsection{Limitations}

\textbf{Branch cut restriction.}
Theorem~\ref{thm:conj} requires $|\mathrm{Im}(z)| < \pi$ for
the principal branch of the logarithm. All data in this work
satisfies this condition; extensions to $|\mathrm{Im}(z)| \geq \pi$
would require branch-tracking.

\textbf{Real-valued data.}
$\emlstar$ is structurally identical to $\eml$ on $\mathbb{R}$.
The method is only informative on data with genuine complex structure.
This is a mathematical fact, not a limitation of the implementation.

\textbf{Neural world model.}
We attempted to build a JEPA-style neural network (complex-valued,
with separate holomorphic and anti-holomorphic pathways).
In a supervised setting, the model achieves 88\% classification
accuracy, but without supervision, the learned gate does not
differentiate the two pathways.
The root cause is that $|\conj{z}| = |z|$: a gate that receives
only magnitudes cannot distinguish holomorphic from anti-holomorphic
dynamics. This remains an open problem.

\subsection{Coordinate representations}

We tested $\emlstar$ on the Schwarzschild metric in two formulations:
the standard Hilbert form (which crosses the horizon $r = 2M$ and
produces complex values) yields $\emlstar = \mathrm{True}$,
while the original Schwarzschild form (which avoids the interior)
yields $\emlstar = \mathrm{False}$.
This is a result about coordinate representations, not about
physical reality: the coordinate singularity at $r = 2M$ has been
understood since Kruskal (1960).


% =====================================================================
\section{Conclusion}
\label{sec:conclusion}

We have shown that the $\emlstar$ operator, combined with symbolic
regression, provides an automatic, interpretable criterion for
detecting anti-holomorphic structure in complex dynamical systems.
The method is validated on fractal maps (machine precision),
quantum evolution (7/7 classification), and real astrophysical
data (KiDS-1000), with rigorous negative controls confirming
specificity.

The $\emlstar$ framework opens several directions:
(i) application to natively complex astrophysical data
(gravitational lensing shear fields, interferometric visibilities);
(ii) integration as a structural prior in neural world models;
(iii) extension to quaternionic or Clifford-algebraic settings.

All code and data are available at
\url{https://github.com/antparis/eml_star}.


% =====================================================================
\section*{Acknowledgments}

The author thanks A.~Odrzywołek for the EML operator framework,
and the KiDS collaboration for public data access.
Symbolic regression was performed with
PySR~\citep{cranmer2023pysr}.
This work was conducted independently without institutional
affiliation or funding.

\bibliographystyle{plainnat}
\begin{thebibliography}{9}

\bibitem[Cranmer(2023)]{cranmer2023pysr}
Cranmer, M. (2023).
Interpretable machine learning for science with PySR and SymbolicRegression.jl.
\textit{arXiv:2305.01582}.

\bibitem[Giblin et~al.(2021)]{giblin2021kids1000}
Giblin, B., et~al. (2021).
KiDS-1000 catalogue: Weak gravitational lensing shear measurements.
\textit{Astronomy \& Astrophysics}, 645, A105.


\bibitem[Petrillo et~al.(2019)]{petrillo2019linKS}
Petrillo, C.~E., et~al. (2019).
LinKS: discovering galaxy-scale strong lenses in the Kilo-Degree Survey
using convolutional neural networks.
\textit{Monthly Notices of the Royal Astronomical Society}, 484(3), 3879--3896.

\bibitem[Li et~al.(2020)]{li2020kids}
Li, R., Napolitano, N.~R., Tortora, C., et~al. (2020).
New high-quality strong lens candidates with deep learning in the Kilo-Degree Survey.
\textit{The Astrophysical Journal}, 899(1), 30.

\bibitem[Li et~al.(2021)]{li2021kids}
Li, R., Napolitano, N.~R., Spiniello, C., et~al. (2021).
High-quality strong lens candidates in the final Kilo-Degree Survey footprint.
\textit{The Astrophysical Journal}, 923(1), 16.

\bibitem[Odrzywołek(2026)]{odrzywołek2026eml}
Odrzywołek, A. (2026).
All elementary functions from a single binary operator.
\textit{arXiv:2603.21852}.

\end{thebibliography}

\end{document}
